% Nouveau EGA V (cote 25): archive cote 25, site page 180 (PDF page 181).
% typescript. Draft compared: none.
% AI-transcribed by Claude Opus 5.5 (claude-opus-5-5, adaptive thinking, effort high), Message Batches API, batch ega25b (msgbatch_01A7n5YuXubUPNK1zsepi8RG), 2026-11-13; not yet checked by eye.
\scanpage{25/180}
\mspage{7}

fidèlement plat sur S, ce qui prouvera  (ii)$\Longrightarrow$(iii bis) dans  16.9.

Démonstration de 16.11.  Comme les deux membres sont lisses sur S,
il suffit de prouver la platitude fibre par fibre, ce qui nous ramène au
cas où S est le spectre d'un corps alg clos k.  On peut \xout{\illegible{1}} supposer
évidemment $\underline{E}_o=\underline{\unsure{B|E}}$, donc $(**)$ prend la forme
$(***)$   $Gl(\underline{E})\longrightarrow\underline{Q}(\underline{E})^*$ ,
ce morphisme étant déduit des opérations naturelles du schéma en groupes
$Gl(\underline{E})$ sur $\underline{Q}(\underline{E})^*$ ,par $u\rightsquigarrow u(q_o)$, où $q_o$ est une section de $\underline{Q}(\underline{E})^*$ corres-%
pondant à la forme standart $Q_o$ . Or par le théorème de platitude \unsure{générique|génériaque}\hand{,}
comme $\underline{Q}(\underline{E})^*$ est lisse sur k,donc réduit, il existe un ouvert dense U de
$\underline{Q}(\underline{E})^*$ au dessus duquel le morphisme précédent soit plat.  Pour tout \margin{\unsure{dégager un lemme sur les morphismes d'espaces homogènes}}
\xout{$u\in Gl(E)(k)$}$u\in Gl(\underline{E})(k)$, i.e. pour tout automorphisme u de $\underline{E}$, u(U) satis-%
fera alors à la même condition, et il suffit\hand{,}\inserted{pour établir la platitude} de prouver\hand{),} \xout{\illegible{1}}
de $(***)$\hand{,}que les u(U) recouvrent $\underline{Q}(\underline{E})^*$ . Or pour ceci il suffit de prouver
que $Gl(\underline{E})(k)=$ Aut$(\underline{E})$ opère transitivement dans $\underline{Q}(\underline{E})^*(k)$, ensemble des
formes quadratiques \xout{\illegible{1}} sur \hand{$\checkmark$}E à discriminant non nul, \xout{\illegible{1}} i.e. que
\deleted{(}$***)$ est surjectif, ce \hand{\unsure{|}}qui prouvera alors 16.11.  On est donc ramené à
prouver le

\uline{Lemme} 16.11.1.  Soit Q une forme quadratique en n variables sur un corps k
alg clos, telle que $\overset{\hand{$\sim$}}{d}(Q)\neq 0$ , alors Q est isomorphe à la forme standart
(i.e. se met sous forme standart par un choix convenable de la base).

Lorsque E est de rang pair ou k est de car $\neq 2$ , alors l'hypothèse
signifie que Q est non dégénérée, et la conclusion se trouve dans Bourbaki
Dans le cas contraire (k de car 2, \xout{\unsure{dim}} rang E impair)  on sait que Q est
dégénérée, \xout{\unsure{donc}} soit $E_1$ une droite dans E qui est dans le noyau de la
forme associée, $E_2$ un supplémentaire, de sorte que Q se met sous forme

\margin{180}
